\section{Results}
\label{sec:results}

Unless stated otherwise, results are for the \texttt{enhits} pool
(EnamineHTS, DOCK3.7), with a test set of \num{105215} molecules held out from
\num{2104302} usable rows, under a single random seed.
\todo{Repeat all headline comparisons over five seeds and report dispersion;
single-seed numbers are provisional throughout.}

\subsection{The descriptor baseline and its learning curve}
\label{sec:results-baseline}

Table~\ref{tab:baseline-curve} reports the descriptor baseline as a function of
training set size. Performance saturates early: increasing the training set by
a factor of \num{38}, from \num{50000} to \num{1893872} molecules, improves
Pearson $r$ by \num{0.028}. This is the expected behaviour of a \num{12}-feature
model, which exhausts its capacity long before it exhausts the available data.

\begin{table}[t]
\centering
\caption{Descriptor baseline on \texttt{enhits}, as a function of training set
size. Test set \num{105215} molecules.}
\label{tab:baseline-curve}
\small
\begin{tabular}{rrrrrr}
\toprule
Train size & Pearson $r$ & Spearman $\rho$ & RMSE & $\mathrm{EF}@1\%$ & $\mathrm{EF}@0.1\%$ \\
\midrule
\num{50000}   & 0.8748 & 0.8613 & 0.4611 & 37.5 & 238 \\
\num{200000}  & 0.8802 & 0.8669 & 0.4517 & 38.4 & 248 \\
\num{1000000} & 0.8872 & 0.8744 & 0.4391 & 38.7 & 257 \\
\num{1893872} & 0.8907 & 0.8783 & 0.4326 & 40.7 & 286 \\
\bottomrule
\end{tabular}
\end{table}

The practical consequence is that the bar a learned representation must clear
is not only high but cheap: a Pearson $r$ of \num{0.88} is available from
twelve descriptors and \num{200000} docked molecules, with no embedding
storage and no pretraining.

\subsection{The baseline is not an artifact of molecular size}
\label{sec:results-size}

Heavy-atom count alone correlates $-0.659$ with the DOCK3.7 score on this pool,
which raises the possibility that the baseline's apparent strength is simply
size encoding. Table~\ref{tab:sizecontrol} shows that it is not.

\begin{table}[t]
\centering
\caption{Size-controlled evaluation on \texttt{enhits}. ``Raw'' is the
conventional metric; ``partial'' and ``residual'' remove a linear and a
degree-3 dependence on heavy-atom count respectively; ``within-stratum''
quantities are computed inside ten equal-count heavy-atom bins and weighted by
bin size. The first row is the degenerate size-only predictor.}
\label{tab:sizecontrol}
\small
\begin{tabular}{lrrrrr}
\toprule
 & \multicolumn{3}{c}{Pearson $r$} & \multicolumn{2}{c}{Within-stratum} \\
\cmidrule(lr){2-4}\cmidrule(lr){5-6}
Model & raw & partial & residual & $r$ & $\mathrm{EF}@1\%$ \\
\midrule
Size only ($-$heavy atoms) & 0.6592 & 0.000 & 0.000 & 0.1478 & 2.6 \\
\midrule
Descriptors, \num{50000}   & 0.8748 & 0.7647 & 0.7598 & 0.7692 & 25.8 \\
Descriptors, \num{200000}  & 0.8802 & 0.7756 & 0.7710 & 0.7797 & 28.1 \\
Descriptors, \num{1000000} & 0.8872 & 0.7896 & 0.7853 & 0.7933 & 29.6 \\
Descriptors, \num{1893872} & 0.8907 & 0.7966 & 0.7924 & 0.8000 & 30.1 \\
\bottomrule
\end{tabular}
\end{table}

Two observations follow. First, the control behaves as intended: the size-only
predictor attains a substantial raw correlation of \num{0.659} but collapses to
a within-stratum correlation of \num{0.148} and an enrichment factor of
\num{2.6}, close to the chance value of \num{1.0}. A feature that is
near-constant within its own stratum cannot discriminate inside it.

Second, the descriptor baseline loses comparatively little. At the largest
training size its correlation falls from \num{0.891} to \num{0.797} under
partial correlation and \num{0.800} within strata, and its top-\SI{1}{\percent}
enrichment from \num{40.7} to \num{30.1}. Approximately \SI{26}{\percent} of
the raw enrichment is attributable to size; the remainder reflects genuine
chemical information carried by the non-size descriptors.

This result is important for the interpretation of everything that follows. A
descriptor baseline that was merely a size proxy would be a weak and somewhat
unfair comparator. One that retains a within-stratum enrichment of \num{30}
is a substantive target, and failure to exceed it is correspondingly
informative.

\subsection{Learned representations}
\label{sec:results-backbones}

\todo{Primary results table pending completion of the backbone sweep. Populate
with: backbone, dimensionality, storage, raw and size-controlled metrics, at
each of four training budgets. Expected structure below.}

\begin{table}[t]
\centering
\caption{Learned representations versus the descriptor baseline on
\texttt{enhits} at the largest training budget. \todo{pending}}
\label{tab:backbones-results}
\small
\begin{tabular}{lrrrrrr}
\toprule
Representation & Dim. & Pearson $r$ & partial $r$ & within $r$ & $\mathrm{EF}@1\%$ & within $\mathrm{EF}@1\%$ \\
\midrule
Descriptors & 12 & 0.8907 & 0.7966 & 0.8000 & 40.7 & 30.1 \\
\midrule
MHG-GED     & 1024 & \pending & \pending & \pending & \pending & \pending \\
GROVER      & 3400 & \pending & \pending & \pending & \pending & \pending \\
Uni-Mol2    &  768 & \pending & \pending & \pending & \pending & \pending \\
MolFormer   &  768 & \pending & \pending & \pending & \pending & \pending \\
SMI-TED     &  768 & \pending & \pending & \pending & \pending & \pending \\
\bottomrule
\end{tabular}
\end{table}

\todo{Learning-curve figure: Pearson $r$ against training size, one line per
representation, with the descriptor baseline drawn as a horizontal reference.
The scientifically informative feature is curve \emph{shape}: a representation
that begins below the baseline but continues to improve where the baseline has
flattened has capacity that additional docking would unlock, which is a
different conclusion from one that flattens at the same point.}

\subsection{Robustness to the choice of prediction head}
\label{sec:results-robust}

A fixed prediction head invites the objection that any observed deficit of
learned representations reflects an architecture tuned, in effect, for
low-dimensional descriptors. We therefore repeat the comparison with two
additional model families on the same representations and splits: ridge
regression, and gradient-boosted decision trees.

\todo{Run and report. If the ordering between representations is preserved
across three model families of markedly different inductive bias, the result
concerns the representation rather than the head. This is the single most
important robustness check in the paper.}

\subsection{Late fusion}
\label{sec:results-fusion}

\todo{Rerun fusion on converged models; the numbers below are from an
undertrained preliminary sweep and are reported only to indicate the direction
of the effect.}

In a preliminary experiment, averaging the predictions of a strong
representation with two weaker ones reduced Pearson $r$ from \num{0.938} to
\num{0.839}. Prediction averaging weights all members equally, so the
combination is pulled toward its weakest constituent. Prior reports that late
fusion reliably improves on single
modalities~\citep{martinezgalindo2026multimodality} were obtained in a regime
where member representations performed comparably; that condition does not hold
here, and unweighted fusion is correspondingly harmful.

\todo{Quantify under the converged sweep, and compare unweighted averaging
against rank averaging and validation-weighted averaging. A gating rule ---
fuse only members that clear the descriptor baseline --- is the natural
remedy and should be evaluated.}

\subsection{Transfer across targets}

\todo{Repeat the headline comparison on \texttt{ampc} and \texttt{d4}. Both
use DOCK3.7, so agreement across them establishes robustness to target but not
to scoring function. Claims about docking in general require the Glide pool,
which is presently unavailable (Section~\ref{sec:data}).}
